\documentclass[aspectratio=169]{beamer}
\usepackage[utf8]{inputenc}
\usepackage[T1]{fontenc}
\usepackage[es-ES]{babel}
\usepackage{graphicx}
\usepackage{tikz}
\usepackage{listings}
\usepackage{xcolor}
\usepackage{amsmath}

% Tema
\usetheme{Madrid}
\usecolortheme{default}

% Colores personalizados
\definecolor{unsazul}{RGB}{30, 60, 120}
\definecolor{unsverde}{RGB}{0, 100, 0}

\setbeamercolor{structure}{fg=unsazul}
\setbeamercolor{palette primary}{bg=unsazul,fg=white}
\setbeamercolor{palette secondary}{bg=unsazul!80,fg=white}
\setbeamercolor{palette tertiary}{bg=unsazul!60,fg=white}

% Configuración de listings
\lstset{
    basicstyle=\ttfamily\small,
    keywordstyle=\color{unsazul}\bfseries,
    commentstyle=\color{gray}\itshape,
    stringstyle=\color{unsverde},
    showstringspaces=false,
    breaklines=true,
    frame=single,
    backgroundcolor=\color{gray!10}
}

% Información del documento
\title{Métodos de Análisis de Datos 1}
\subtitle{Clase 6: Normalización e Integración}
\author{Universidad Nacional del Sur}
\date{}

\begin{document}

% Portada
\begin{frame}
\titlepage
\end{frame}

% Contenidos
\begin{frame}{Contenidos de hoy}
\tableofcontents
\end{frame}

\section{Normalización y Estandarización}

\begin{frame}{El problema de las escalas}
\begin{exampleblock}{Ejemplo}
\begin{center}
\begin{tabular}{|c|c|}
\hline
\textbf{Edad} & \textbf{Ingreso} \\
\hline
25 & 50000 \\
30 & 65000 \\
45 & 80000 \\
\hline
\end{tabular}
\end{center}

El ingreso tiene valores mucho mayores que la edad
\end{exampleblock}

\pause

\begin{alertblock}{Problemas}
\begin{itemize}
    \item Algoritmos sensibles a escala (KNN, SVM, redes neuronales)
    \item Variables con mayor rango dominan el modelo
    \item Convergencia lenta en algoritmos de optimización
    \item Interpretación difícil de coeficientes
\end{itemize}
\end{alertblock}
\end{frame}

\begin{frame}{Solución: Escalar las variables}
\begin{block}{¿Qué es escalar?}
Transformar variables a rangos comparables
\end{block}

\begin{block}{Dos técnicas principales}
\begin{enumerate}
    \item \textbf{Normalización (Min-Max):} Escalar a rango [0, 1]
    \item \textbf{Estandarización (Z-score):} Transformar a media 0 y desviación estándar 1
\end{enumerate}
\end{block}

\pause

\begin{alertblock}{¿Cuándo escalar?}
\begin{itemize}
    \item ✓ KNN, SVM, redes neuronales
    \item ✓ Algoritmos de clustering (K-means)
    \item ✓ Regresión regularizada (Ridge, LASSO)
    \item ✗ Árboles de decisión, Random Forest
\end{itemize}
\end{alertblock}
\end{frame}

\begin{frame}{Min-Max Scaling (Normalización)}
\begin{block}{Fórmula}
$$X_{norm} = \frac{X - X_{min}}{X_{max} - X_{min}}$$

Transforma valores al rango [0, 1]
\end{block}

\begin{exampleblock}{Ejemplo}
Variable: [10, 20, 30, 40, 50]

\begin{align*}
X_{norm}(10) &= \frac{10 - 10}{50 - 10} = 0.00 \\
X_{norm}(30) &= \frac{30 - 10}{50 - 10} = 0.50 \\
X_{norm}(50) &= \frac{50 - 10}{50 - 10} = 1.00
\end{align*}

Resultado: [0.00, 0.25, 0.50, 0.75, 1.00]
\end{exampleblock}
\end{frame}

\begin{frame}[fragile]{Min-Max: Código}
\begin{columns}[t]
\column{0.5\textwidth}
\textbf{Python:}
\begin{lstlisting}[language=Python]
from sklearn.preprocessing import MinMaxScaler

# Normalizar
scaler = MinMaxScaler()
df[['edad', 'ingreso']] = scaler.fit_transform(
    df[['edad', 'ingreso']]
)

# Rango personalizado [-1, 1]
scaler = MinMaxScaler(
    feature_range=(-1, 1)
)
df[['edad', 'ingreso']] = scaler.fit_transform(
    df[['edad', 'ingreso']]
)
\end{lstlisting}

\column{0.5\textwidth}
\textbf{R:}
\begin{lstlisting}[language=R]
# Manual
normalizar <- function(x) {
  (x - min(x, na.rm = TRUE)) / 
  (max(x, na.rm = TRUE) - min(x, na.rm = TRUE))
}

df <- df %>%
  mutate(across(
    c(edad, ingreso), 
    normalizar
  ))

# Con caret
library(caret)
preproc <- preProcess(
  df[, c('edad', 'ingreso')], 
  method = "range"
)
df[, c('edad', 'ingreso')] <- predict(preproc, df)
\end{lstlisting}
\end{columns}
\end{frame}

\begin{frame}{Standardization (Z-score)}
\begin{block}{Fórmula}
$$X_{std} = \frac{X - \mu}{\sigma}$$

Transforma a media 0 y desviación estándar 1
\end{block}

\begin{exampleblock}{Ejemplo}
Variable: [10, 20, 30, 40, 50]\\
Media ($\mu$) = 30, Desv. Est. ($\sigma$) = 14.14

\begin{align*}
X_{std}(10) &= \frac{10 - 30}{14.14} = -1.41 \\
X_{std}(30) &= \frac{30 - 30}{14.14} = 0.00 \\
X_{std}(50) &= \frac{50 - 30}{14.14} = 1.41
\end{align*}

Resultado: [-1.41, -0.71, 0.00, 0.71, 1.41]
\end{exampleblock}
\end{frame}

\begin{frame}[fragile]{Standardization: Código}
\begin{columns}[t]
\column{0.5\textwidth}
\textbf{Python:}
\begin{lstlisting}[language=Python]
from sklearn.preprocessing import StandardScaler

scaler = StandardScaler()
df[['edad', 'ingreso']] = scaler.fit_transform(
    df[['edad', 'ingreso']]
)

# Verificar
print(df['edad'].mean())  # ~0
print(df['edad'].std())   # ~1
\end{lstlisting}

\column{0.5\textwidth}
\textbf{R:}
\begin{lstlisting}[language=R]
# Estandarizar
df <- df %>%
  mutate(across(
    c(edad, ingreso), 
    scale
  ))

# Con caret
preproc <- preProcess(
  df[, c('edad', 'ingreso')], 
  method = c("center", "scale")
)
df[, c('edad', 'ingreso')] <- predict(preproc, df)
\end{lstlisting}
\end{columns}
\end{frame}

\begin{frame}{¿Min-Max o Standardization?}
\begin{columns}[t]
\column{0.5\textwidth}
\textbf{Min-Max:}
\begin{itemize}
    \item ✓ Rango específico necesario
    \item ✓ Redes neuronales
    \item ✓ Algoritmos de imágenes
    \item ✓ No hay outliers extremos
    \item ✗ Sensible a outliers
\end{itemize}

\column{0.5\textwidth}
\textbf{Standardization:}
\begin{itemize}
    \item ✓ Distribución aproximadamente normal
    \item ✓ Regresión logística, SVM
    \item ✓ PCA, clustering
    \item ✓ Más robusto a outliers
    \item ✗ No acota valores
\end{itemize}
\end{columns}

\vspace{1em}

\begin{exampleblock}{Regla práctica}
\textbf{¿En duda?} → Usar Standardization
\end{exampleblock}
\end{frame}

\begin{frame}{Importante: fit vs transform}
\begin{alertblock}{Error común: Data leakage}
Nunca usar estadísticas del conjunto de test para escalar
\end{alertblock}

\begin{block}{Flujo correcto}
\begin{enumerate}
    \item Dividir en train/test
    \item \texttt{fit()} con datos de train (calcular min/max o media/std)
    \item \texttt{transform()} train con parámetros calculados
    \item \texttt{transform()} test con los MISMOS parámetros
\end{enumerate}
\end{block}

\pause

\begin{exampleblock}{Código correcto}
\texttt{scaler.fit(X\_train)}\\
\texttt{X\_train\_scaled = scaler.transform(X\_train)}\\
\texttt{X\_test\_scaled = scaler.transform(X\_test)}
\end{exampleblock}
\end{frame}

\section{Discretización}

\begin{frame}{Discretización (Binning)}
\begin{block}{¿Qué es?}
Convertir variables continuas en categóricas dividiendo en intervalos
\end{block}

\begin{exampleblock}{Ejemplo}
Edad continua → Grupos de edad

\begin{center}
\begin{tabular}{|c|c|}
\hline
\textbf{Edad} & \textbf{Grupo} \\
\hline
15 & menor \\
25 & joven \\
35 & adulto \\
55 & maduro \\
75 & mayor \\
\hline
\end{tabular}
\end{center}
\end{exampleblock}
\end{frame}

\begin{frame}{¿Por qué discretizar?}
\begin{block}{Ventajas}
\begin{itemize}
    \item Captura relaciones no lineales
    \item Más robusto a outliers
    \item Facilita interpretación
    \item Útil para análisis exploratorio
    \item Algunos modelos lo requieren
\end{itemize}
\end{block}

\pause

\begin{alertblock}{Desventajas}
\begin{itemize}
    \item Pérdida de información
    \item Introduce discontinuidades artificiales
    \item Elección de bins es subjetiva
    \item Puede empeorar performance
\end{itemize}
\end{alertblock}
\end{frame}

\begin{frame}[fragile]{Tipos de discretización}
\begin{block}{1. Bins de igual ancho}
Dividir el rango en intervalos iguales
\end{block}

\begin{lstlisting}[language=Python]
# Python
df['edad_grupo'] = pd.cut(
    df['edad'],
    bins=5,  # 5 intervalos iguales
    labels=['muy_joven', 'joven', 'adulto', 'mayor', 'muy_mayor']
)
\end{lstlisting}

\begin{lstlisting}[language=R]
# R
df <- df %>%
  mutate(edad_grupo = cut(
    edad,
    breaks = 5,
    labels = c('muy_joven', 'joven', 'adulto', 'mayor', 'muy_mayor')
  ))
\end{lstlisting}
\end{frame}

\begin{frame}[fragile]{Bins con puntos específicos}
\begin{lstlisting}[language=Python]
# Python - Bins personalizados
df['edad_grupo'] = pd.cut(
    df['edad'],
    bins=[0, 18, 30, 50, 65, 100],
    labels=['menor', 'joven', 'adulto', 'maduro', 'mayor']
)
\end{lstlisting}

\begin{lstlisting}[language=R]
# R - Bins personalizados
df <- df %>%
  mutate(edad_grupo = cut(
    edad,
    breaks = c(0, 18, 30, 50, 65, 100),
    labels = c('menor', 'joven', 'adulto', 'maduro', 'mayor')
  ))
\end{lstlisting}
\end{frame}

\begin{frame}[fragile]{Bins de igual frecuencia (Quantiles)}
\begin{block}{Concepto}
Cada bin contiene aproximadamente el mismo número de observaciones
\end{block}

\begin{lstlisting}[language=Python]
# Python - Cuartiles
df['ingreso_grupo'] = pd.qcut(
    df['ingreso'],
    q=4,  # 4 cuartiles
    labels=['bajo', 'medio_bajo', 'medio_alto', 'alto']
)
\end{lstlisting}

\begin{lstlisting}[language=R]
# R - Cuartiles
df <- df %>%
  mutate(ingreso_grupo = cut(
    ingreso,
    breaks = quantile(ingreso, probs = seq(0, 1, 0.25), na.rm = TRUE),
    labels = c('bajo', 'medio_bajo', 'medio_alto', 'alto'),
    include.lowest = TRUE
  ))
\end{lstlisting}
\end{frame}

\section{Integración de datos}

\begin{frame}{Combinar múltiples fuentes}
\begin{block}{¿Por qué integrar datos?}
\begin{itemize}
    \item Datos dispersos en múltiples tablas/archivos
    \item Enriquecer con información externa
    \item Análisis requiere combinar fuentes
\end{itemize}
\end{block}

\pause

\begin{block}{Dos operaciones principales}
\begin{enumerate}
    \item \textbf{Joins:} Combinar por columnas comunes (como SQL)
    \item \textbf{Concatenación:} Apilar tablas vertical u horizontalmente
\end{enumerate}
\end{block}
\end{frame}

\begin{frame}{Tipos de Joins}
\begin{center}
\begin{tikzpicture}[scale=0.8]
    % Inner join
    \node at (0,3) {\textbf{Inner Join}};
    \draw[fill=unsazul!30] (0,2) circle (0.8cm);
    \draw[fill=unsverde!30] (1.2,2) circle (0.8cm);
    \draw[fill=unsazul!50!unsverde!50] (0.6,2) circle (0.4cm);
    
    % Left join
    \node at (4,3) {\textbf{Left Join}};
    \draw[fill=unsazul!50] (4,2) circle (0.8cm);
    \draw[fill=unsverde!30] (5.2,2) circle (0.8cm);
    
    % Right join
    \node at (8,3) {\textbf{Right Join}};
    \draw[fill=unsazul!30] (8,2) circle (0.8cm);
    \draw[fill=unsverde!50] (9.2,2) circle (0.8cm);
    
    % Outer join
    \node at (12,3) {\textbf{Outer Join}};
    \draw[fill=unsazul!50] (12,2) circle (0.8cm);
    \draw[fill=unsverde!50] (13.2,2) circle (0.8cm);
\end{tikzpicture}
\end{center}

\begin{itemize}
    \item \textbf{Inner:} Solo filas con match en ambas
    \item \textbf{Left:} Todas de izquierda + matches de derecha
    \item \textbf{Right:} Todas de derecha + matches de izquierda
    \item \textbf{Outer:} Todas las filas de ambas
\end{itemize}
\end{frame}

\begin{frame}[fragile]{Joins en Python}
\begin{lstlisting}[language=Python]
import pandas as pd

# Inner join
df_merged = pd.merge(df1, df2, on='id', how='inner')

# Left join
df_merged = pd.merge(df1, df2, on='id', how='left')

# Join con nombres diferentes
df_merged = pd.merge(
    df1, df2,
    left_on='cliente_id',
    right_on='id',
    how='left'
)

# Join con múltiples columnas
df_merged = pd.merge(
    df1, df2, 
    on=['id', 'fecha'], 
    how='inner'
)
\end{lstlisting}
\end{frame}

\begin{frame}[fragile]{Joins en R}
\begin{lstlisting}[language=R]
library(dplyr)

# Inner join
df_merged <- inner_join(df1, df2, by = "id")

# Left join
df_merged <- left_join(df1, df2, by = "id")

# Join con nombres diferentes
df_merged <- left_join(
  df1, df2,
  by = c("cliente_id" = "id")
)

# Join con múltiples columnas
df_merged <- inner_join(
  df1, df2, 
  by = c("id", "fecha")
)
\end{lstlisting}
\end{frame}

\begin{frame}[fragile]{Concatenación}
\begin{columns}[t]
\column{0.5\textwidth}
\textbf{Python:}
\begin{lstlisting}[language=Python]
# Vertical (apilar filas)
df_combined = pd.concat(
    [df1, df2], 
    axis=0,           # 0 = vertical
    ignore_index=True # Resetear índice
)

# Horizontal (añadir columnas)
df_combined = pd.concat(
    [df1, df2], 
    axis=1            # 1 = horizontal
)
\end{lstlisting}

\column{0.5\textwidth}
\textbf{R:}
\begin{lstlisting}[language=R]
library(dplyr)

# Vertical
df_combined <- bind_rows(
  df1, df2
)

# Horizontal
df_combined <- bind_cols(
  df1, df2
)
\end{lstlisting}
\end{columns}

\vspace{1em}

\begin{alertblock}{Cuidado}
Al concatenar horizontalmente, asegurar que las filas corresponden
\end{alertblock}
\end{frame}

\section{Resumen}

\begin{frame}{Resumen de la clase}
\begin{itemize}
    \item \textbf{Normalización (Min-Max):} Escalar a [0, 1], sensible a outliers
    \item \textbf{Estandarización (Z-score):} Media 0, std 1, más robusto
    \item Usar \texttt{fit()} solo con train, \texttt{transform()} con train y test
    \item \textbf{Discretización:} Convertir continuas en categóricas
    \item Bins: igual ancho, personalizados, igual frecuencia
    \item \textbf{Joins:} Combinar tablas por columnas comunes
    \item Tipos: inner, left, right, outer
    \item \textbf{Concatenación:} Apilar vertical u horizontal
\end{itemize}

\vspace{1em}

\begin{block}{Próxima clase}
\textbf{APIs REST, web scraping y tidy data}
\end{block}
\end{frame}

\begin{frame}{Tarea}
\begin{enumerate}
    \item Usar el dataset transformado de la clase anterior
    \item Normalización/Estandarización:
    \begin{itemize}
        \item Aplicar Min-Max a variables numéricas
        \item Aplicar Standardization a las mismas variables
        \item Comparar distribuciones antes/después
    \end{itemize}
    \item Discretización:
    \begin{itemize}
        \item Elegir 2 variables continuas
        \item Crear bins con diferentes estrategias
        \item Visualizar resultados
    \end{itemize}
    \item Si tenés múltiples tablas, practicar joins
    \item Documentar todo en el notebook
    \item Subir a GitHub
\end{enumerate}
\end{frame}

\begin{frame}[standout]
\Huge ¿Preguntas?
\end{frame}

\end{document}
